\documentclass[11pt]{article}
\usepackage[margin=0.75in]{geometry}
\usepackage{fontspec}
\setmainfont{Times New Roman}
\newfontfamily\cjkfont{AppleMyungjo}
\newcommand{\ko}[1]{{\cjkfont #1}}
\usepackage{enumitem}
\usepackage{setspace}
\singlespacing
\setlength{\parindent}{0pt}
\setlength{\parskip}{4pt}
\pagestyle{empty}

\begin{document}

\begin{center}
    {\large \textbf{Proxy Voting for Children}}\\[3pt]
    Daniel Jeong \quad|\quad PHIL 312
\end{center}

\vspace{-3pt}
\rule{\textwidth}{0.4pt}

Children have no vote, yet they bear the consequences for longer than anyone voting today. Parents already act on their children's behalf in every other domain: medical decisions, schooling, legal matters, where they live. Voting is the exception.

\textbf{The Proposal:} Each minor child generates one proxy vote, cast by their legal guardians---half a vote each where two exist. The vote expires when the child turns 18 and casts their own.

\textbf{Why Parents Make Good Proxies.}

\textit{Entanglement.} The proxy bears the cost when things go wrong for the principal. A parent who votes against their child's interest goes home to the consequences. When schools fail, the parent watches their child struggle through them. When healthcare fails, the parent absorbs the bills. When the labor market fails, an adult child returns home. Forfeiting custody carries such legal and social cost that it is not a real option.

\textit{Alignment.} Parents are not fully aligned with their children's interests. Full alignment would imply total self-sacrifice, which most parents do not undertake. But they are reliably aligned in one direction: a parent has a direct interest in preventing their child's failure---in keeping them from becoming financially, physically, or emotionally dependent indefinitely. Children share this interest. The alignment is imperfect, but it runs in the right direction on the questions a proxy must vote on.

\textbf{If Proxies Work, Why Not Epistocracy?}

If we trust parents to vote on behalf of their children, why not allow experts to vote on behalf of everyone? Brennan argues exactly this: the most knowledgeable should rule, not the public. The same two criteria apply.

\textit{Entanglement.} Experts are not entangled with those they would govern. A Nobel laureate's life does not depend on what happens to high-school dropouts. When ordinary lives are disrupted by bad policy, the expert's life is relatively unscathed.

\textit{Alignment.} The relationship runs in the opposite direction. Credentialed self-interest runs \textit{against} the broader population, not toward it. Medical associations limit how many doctors may practice. Bar associations restrict who may give legal advice. Every credentialed group makes itself smaller and more exclusive at the expense of outsiders. Where parental selfishness happens to serve the child, expert selfishness actively harms those outside the credential. Education does not eliminate this; it merely makes the defense more articulate.

So epistocracy fails both criteria. Parents pass both.

\textbf{Objections.}

\begin{itemize}[nosep, leftmargin=1.5em, itemsep=2pt]
    \item \textit{The proxy votes on every matter, not just family policy.} This is true of every voter. Childless adults vote on schools. City dwellers vote on rural policy. Retirees vote on student loans. The proxy follows the same rule, and there is no principled reason it should not.
    \item \textit{Why not lower the voting age instead?} Lowering the voting age is compatible with the proposal, not competitive with it. If the age were lowered to 14, the proxy would simply cover ages 0--13 rather than 0--17. The conflict arises only if one proposes to enfranchise infants directly, which dissolves the rationale for any voting age.
\end{itemize}

\textbf{Core Claim:} We already entrust parents with medical, legal, financial, and educational decisions on behalf of their children. Voting alone remains the exception. There is no principled reason it should remain so.

\end{document}